% !TEX root = ../matan.tex

\section{Формула для коэффициентов разложения в ряд аналитической функции. Несовпадение классов бесконечно дифференцируемых и аналитических функций.}

\begin{Theorem}{Формула для коэффициентов (Тейлора)}{}
	Пусть функция $f$ в окрестности точки $x_0$ представима степенным рядом
	\[
	f(x) = \sum_{n=0}^{\infty} a_n (x - x_0)^n, \qquad |x - x_0| < R,\ R > 0.
	\]
	Тогда $f$ бесконечно дифференцируема в этой окрестности и её коэффициенты однозначно определяются формулой
	\[
	a_k = \frac{f^{(k)}(x_0)}{k!}, \qquad k = 0, 1, 2, \ldots
	\]
\end{Theorem}

\begin{proof}
	По теореме о почленном дифференцировании степенного ряда (вопрос~32) ряд можно дифференцировать любое число раз внутри интервала сходимости, причём радиус не меняется. Дифференцируя $k$ раз,
	\[
	f^{(k)}(x) = \sum_{n \ge k} a_n\, n(n-1)\cdots(n - k + 1)\,(x - x_0)^{n-k}.
	\]
	Подставляя $x = x_0$, видим, что выживает лишь член $n = k$ (остальные содержат положительную степень $(x - x_0)$): $f^{(k)}(x_0) = a_k \cdot k!$. Отсюда $a_k = f^{(k)}(x_0)/k!$.
\end{proof}

\begin{Corollary}{Единственность разложения}{}
	Если функция представима степенным рядом с центром $x_0$, то этот ряд единствен: его коэффициенты обязаны совпадать с коэффициентами Тейлора $f^{(k)}(x_0)/k!$. В частности, разложение аналитической функции в степенной ряд в данной точке определено однозначно.
\end{Corollary}

\begin{Remark}{Следствие для аналитических функций}{}
	Таким образом, если $f$ аналитична в $x_0$, то её степенной ряд с необходимостью есть ряд Тейлора $\sum_{k} \dfrac{f^{(k)}(x_0)}{k!}(x - x_0)^k$, и он сходится к $f$ в окрестности $x_0$.
\end{Remark}

\subsection*{$C^\infty \neq$ аналитичность}

Бесконечная дифференцируемость \emph{не} влечёт аналитичность: у функции класса $C^\infty$ ряд Тейлора может сходиться, но не к самой функции.

\begin{Example}{Гладкая, но не аналитическая функция}{}
	Рассмотрим
	\[
	f(x) = \begin{cases} e^{-1/x^2}, & x \neq 0, \\ 0, & x = 0. \end{cases}
	\]
	Тогда $f \in C^\infty(\RR)$ и $f^{(k)}(0) = 0$ для всех $k$, поэтому ряд Тейлора в нуле тождественно равен нулю, тогда как $f(x) \neq 0$ при $x \neq 0$. Значит, ряд Тейлора не сходится к $f$ ни в какой проколотой окрестности нуля, и $f$ не аналитична в $0$ (хотя и бесконечно дифференцируема).
\end{Example}

\begin{Statement}{Обоснование примера}{}
	Для функции $f$ из примера справедливо $f \in C^\infty(\RR)$ и $f^{(k)}(0) = 0$ при всех $k$. 
\end{Statement}

\begin{proof}
	\textbf{Быстрое убывание.} Для любого $m \ge 0$
	\[
	\lim_{x \to 0} \frac{e^{-1/x^2}}{x^m} = 0.
	\]
	Действительно, при $x \to 0$ положим $t = 1/x^2 \to +\infty$; тогда $|x|^{-m} = t^{m/2}$, и $t^{m/2} e^{-t} \to 0$, так как экспонента растёт быстрее любой степени.

	\textbf{Вид производных при $x \neq 0$.} По индукции $f^{(k)}(x) = P_k(1/x)\,e^{-1/x^2}$ для $x \neq 0$, где $P_k$~--- некоторый многочлен. База $k = 0$: $P_0 = 1$. Переход: дифференцируя $P_k(1/x) e^{-1/x^2}$, по правилу произведения снова получаем многочлен от $1/x$, умноженный на $e^{-1/x^2}$ (производная $e^{-1/x^2}$ равна $\tfrac{2}{x^3} e^{-1/x^2}$).

	\textbf{Производные в нуле равны нулю.} Индукция по $k$. При $k = 0$ имеем $f(0) = 0$. Пусть $f^{(k-1)}(0) = 0$. Тогда
	\[
	f^{(k)}(0) = \lim_{x \to 0} \frac{f^{(k-1)}(x) - f^{(k-1)}(0)}{x} = \lim_{x \to 0} \frac{P_{k-1}(1/x)\,e^{-1/x^2}}{x}.
	\]
	Выражение $\dfrac{P_{k-1}(1/x)}{x}$~--- многочлен от $1/x$, то есть конечная сумма членов вида $C\,x^{-m}$; каждый из них, умноженный на $e^{-1/x^2}$, стремится к нулю по первому пункту. Значит, $f^{(k)}(0) = 0$. Тем самым $f$ бесконечно дифференцируема в нуле и все её производные там нулевые.
\end{proof}

\begin{Remark}{}{}
	Итак, $f \in C^\infty(\RR)$, но $f$ не аналитична в нуле: включение классов строгое,
	\[
	\{\text{аналитические функции}\} \subsetneq C^\infty.
	\]
	Над $\C$ ситуация иная: голоморфная (однажды комплексно дифференцируемая) функция автоматически аналитична.
\end{Remark}
